\chapter{Descripción del Sistema}

\section{Arquitectura General}

El sistema de microscopía desarrollado se concibe como una carga útil experimental compacta, modular y de bajo consumo, diseñada para su integración en el proyecto INTISAT. La arquitectura general prioriza la simplicidad funcional, la trazabilidad de los datos y la reproducibilidad experimental, en concordancia con las restricciones propias de plataformas tipo CubeSat.

Desde un punto de vista funcional, el sistema se organiza en cuatro bloques principales: adquisición óptica, iluminación, control y procesamiento, y almacenamiento de datos. Esta descomposición permite aislar las funciones físicas de captura de imagen de las funciones lógicas de control y análisis, facilitando tanto la verificación del sistema como su futura ampliación.

El sistema opera bajo una arquitectura centralizada, en la cual una computadora de carga útil coordina la adquisición de imágenes, el control de la iluminación y el registro de metadatos, asegurando coherencia entre las condiciones experimentales y los datos generados.

\begin{figure}[H]
    \centering
    \includegraphics[width=0.65\textwidth]{fig/Fig_03_01_Arquitectura_General_Payload.png}
    \caption{Arquitectura general del payload de microscopía: bloques funcionales y flujo de datos.}
    \label{fig:arch_general}
\end{figure}

\section{Bloque de adquisición óptica}

El bloque de adquisición óptica está compuesto por un sensor CMOS comercial operando en configuración de microscopía sin lentes. En esta arquitectura, la muestra se coloca en contacto cercano con el plano del sensor, eliminando el uso de elementos ópticos refractivos para la formación de imagen.

La imagen se genera a partir de la interacción directa entre la luz incidente y la muestra, siendo registrada como una distribución espacial de intensidad en la matriz de píxeles del sensor. El bloque permite el control manual de parámetros críticos como el tiempo de exposición y la ganancia, deshabilitando mecanismos automáticos para garantizar consistencia entre capturas consecutivas.

Este enfoque permite una arquitectura óptica simple, robusta y adecuada para validaciones experimentales, aunque introduce limitaciones asociadas a la resolución efectiva y a la sensibilidad al acoplamiento óptico muestra--sensor.

\begin{figure}[H]
    \centering
    \includegraphics[width=0.92\textwidth]{fig/Fig_03_02_Lensless_Setup_CMOS_Muestra.png}
    \caption{Esquema conceptual de microscopía sin lentes: acoplamiento muestra--sensor y registro directo del patrón de intensidad.}
    \label{fig:lensless_setup}
\end{figure}

\section{Bloque de iluminación}

El bloque de iluminación proporciona la fuente de luz necesaria para la observación de la muestra y está implementado mediante un sistema plano y programable. Este bloque permite generar patrones de iluminación controlados en intensidad y geometría, adaptándose tanto a la operación nominal como a modos experimentales.

En condiciones nominales, el sistema emplea iluminación aproximadamente uniforme con el objetivo de maximizar el contraste de la muestra y reducir artefactos asociados a variaciones espaciales de intensidad. De manera complementaria, se dispone de la capacidad de generar patrones estructurados de iluminación para evaluación experimental de técnicas computacionales avanzadas.

El diseño del bloque de iluminación prioriza la estabilidad temporal, la reproducibilidad de los patrones generados y la integración compacta con el resto del sistema.

\begin{figure}[H]
    \centering
    \includegraphics[width=0.95\textwidth]{fig/Fig_03_03_Patrones_Iluminacion_Uniforme_Rejilla_Puntual.png}
    \caption{Modos de iluminación evaluados: (a) uniforme, (b) estructurada tipo rejilla y (c) puntual/dirigida.}
    \label{fig:illum_patterns}
\end{figure}

\section{Bloque de control y procesamiento}

El bloque de control y procesamiento constituye el núcleo lógico del sistema y se encarga de coordinar la operación de todos los subsistemas. Este bloque se implementa mediante una computadora de carga útil que ejecuta el software de control del payload.

Sus funciones principales incluyen la configuración del sensor CMOS, el control del sistema de iluminación, la gestión de los modos operativos y la ejecución de rutinas básicas de calibración, tales como la adquisición de imágenes de referencia para corrección de ruido y no uniformidades.

Asimismo, este bloque gestiona el registro automático de parámetros de adquisición y métricas básicas de calidad de imagen, garantizando la trazabilidad experimental de cada captura. El procesamiento computacional avanzado se concibe como una etapa complementaria y no afecta la operación nominal del sistema.

\subsection{Computadora de carga útil: Raspberry Pi 5}

La computadora de carga útil del sistema de microscopía está implementada mediante una Raspberry Pi 5, seleccionada como plataforma central de control, adquisición y gestión de datos del payload.

La Raspberry Pi 5 actúa como elemento coordinador de todos los subsistemas, encargándose de la configuración y control del sensor CMOS, la gestión del sistema de iluminación, la ejecución de los modos operativos y el almacenamiento de los datos generados. Su capacidad de procesamiento y su compatibilidad con interfaces de alta velocidad permiten una integración directa y estable con el sistema de adquisición óptica.

La elección de la Raspberry Pi 5 responde a criterios de disponibilidad, flexibilidad, bajo consumo relativo y amplia adopción en entornos de instrumentación experimental. Asimismo, su ecosistema de software y soporte para sistemas operativos basados en Linux facilita la implementación de rutinas de control, registro de metadatos y procesamiento básico de imágenes.

Desde el punto de vista arquitectónico, la Raspberry Pi 5 se considera un componente funcional del payload y no una plataforma de propósito general. Su operación se limita estrictamente a las tareas necesarias para la adquisición, control y trazabilidad experimental del sistema de microscopía.

\begin{table}[H]
\centering
\caption{Características técnicas relevantes de la Raspberry Pi 5 como computadora de carga útil}
\begin{tabular}{|l|l|}
\hline
\textbf{Parámetro} & \textbf{Valor relevante} \\
\hline
Modelo & Raspberry Pi 5 \\
Arquitectura & ARM de 64 bits \\
Rol en el sistema & Control y adquisición \\
Interfaces utilizadas & CSI (cámara), I2C (iluminación), GPIO \\
Sistema operativo & Linux embebido \\
Capacidad de procesamiento & Adecuada para control y análisis básico \\
Almacenamiento & Memoria externa (microSD / SSD) \\
\hline
\end{tabular}
\end{table}

\section{Bloque de almacenamiento y datos}

El bloque de almacenamiento se encarga de la persistencia de la información generada durante la operación del sistema. Incluye tanto las imágenes adquiridas como los metadatos asociados a cada captura.

Las imágenes se almacenan en formatos digitales sin compresión o con compresión mínima, preservando la información original capturada por el sensor. De manera complementaria, se generan archivos de metadatos estructurados que contienen parámetros de adquisición, configuraciones de iluminación y marcas temporales.

Esta organización de datos permite una trazabilidad completa del proceso experimental y facilita el análisis posterior en etapas de validación y procesamiento.

\begin{figure}[H]
    \centering
    \includegraphics[width=0.92\textwidth]{fig/Fig_03_04_Estructura_Datos_Imagen_Metadatos.png}
    \caption{Estructura de datos propuesta: imágenes + metadatos para trazabilidad experimental.}
    \label{fig:data_schema}
\end{figure}

\section{Flujo funcional del sistema}

El flujo funcional del sistema durante una operación nominal de captura puede resumirse de la siguiente manera:

\begin{enumerate}
    \item Inicialización del sistema de control y configuración del sensor CMOS.
    \item Configuración del patrón e intensidad del sistema de iluminación.
    \item Verificación de parámetros de adquisición y condiciones de exposición.
    \item Captura de la imagen por parte del sensor CMOS.
    \item Registro de métricas básicas de calidad de imagen.
    \item Almacenamiento de la imagen y generación de metadatos asociados.
\end{enumerate}

Este flujo se mantiene constante para la operación nominal del payload, variando únicamente en modos experimentales el número de capturas y la secuencia de iluminación aplicada. La separación clara entre adquisición y procesamiento garantiza que la imagen base represente fielmente la información física capturada por el sistema.

\begin{figure}[H]
    \centering
    \includegraphics[width=0.95\textwidth]{fig/Fig_03_05_Flujo_Funcional_Captura.png}
    \caption{Flujo funcional nominal: control, iluminación, captura, métricas y almacenamiento.}
    \label{fig:functional_flow}
\end{figure}

\begin{table}[H]
\centering
\caption{Resumen de bloques funcionales del sistema}
\begin{tabular}{|l|l|}
\hline
\textbf{Bloque} & \textbf{Función principal} \\
\hline
Adquisición óptica & Captura directa mediante sensor CMOS \\
Iluminación & Iluminación uniforme y estructurada \\
Control y procesamiento & Coordinación y registro de parámetros \\
Almacenamiento y datos & Persistencia de imágenes y metadatos \\
\hline
\end{tabular}
\end{table}
